% ======================================================================
%  Преамбюл на доклада (XeLaTeX)
% ======================================================================
\usepackage[a4paper,top=2.6cm,bottom=2.6cm,left=2.7cm,right=2.4cm,headheight=15pt]{geometry}
\usepackage{amsmath,amssymb,mathtools,bm}
\usepackage[no-math]{fontspec}
\usepackage{polyglossia}
\setdefaultlanguage{bulgarian}
\setotherlanguage{english}
\setmainfont{PT Serif}
\setsansfont{PT Sans}
\setmonofont{DejaVu Sans Mono}[Scale=0.84]
\newfontfamily\cyrillicfont{PT Serif}
\newfontfamily\cyrillicfontsf{PT Sans}
\newfontfamily\cyrillicfonttt{DejaVu Sans Mono}[Scale=0.84]

\usepackage{xcolor}
\definecolor{navy}{HTML}{1F2A44}
\definecolor{accent}{HTML}{2A78D6}
\definecolor{orange}{HTML}{EB6834}
\definecolor{attack}{HTML}{E34948}
\definecolor{muted}{HTML}{52514E}
\definecolor{rule}{HTML}{C9CDD6}
\definecolor{boxbg}{HTML}{F2F5FA}
\definecolor{rowalt}{HTML}{F5F7FA}

\usepackage{graphicx}
\usepackage{tikz}
\usetikzlibrary{arrows.meta,positioning,calc}
\usepackage{booktabs,tabularx,makecell,multirow,array,longtable}
\usepackage[table]{xcolor}
\usepackage{caption,subcaption}
\captionsetup{font={small,sf},labelfont={bf,color=navy},labelsep=period,justification=justified,singlelinecheck=true,skip=6pt}
\captionsetup[table]{position=top}
\usepackage{siunitx}
\sisetup{output-decimal-marker={,},group-separator={\,},group-minimum-digits=4}
\usepackage{enumitem}
\setlist{itemsep=2pt,topsep=4pt,leftmargin=1.6em}
\usepackage{microtype}\setlength{\emergencystretch}{2.5em}
\usepackage{setspace}
\setstretch{1.12}
\setlength{\parskip}{3pt}
\usepackage{float}
\usepackage{adjustbox}\newcommand{\fitbox}[1]{\adjustbox{max width=\linewidth}{#1}}
\usepackage{placeins}\usepackage{needspace}\usepackage{etoolbox}

% --- непрекъсната номерация на формули, фигури и таблици в целия документ ---
\counterwithout{equation}{chapter}
\counterwithout{figure}{chapter}
\counterwithout{table}{chapter}

% --- заглавия ------------------------------------------------------------
\usepackage{titlesec}
\titleformat{\chapter}[display]
  {\sffamily\color{navy}}
  {\normalsize\bfseries\color{accent}\MakeUppercase{\chaptertitlename}\ \thechapter}
  {4pt}
  {\Huge\bfseries}[\vspace{4pt}{\color{accent}\titlerule[1.2pt]}]
\titleformat{name=\chapter,numberless}[display]
  {\sffamily\color{navy}}{}{0pt}{\Huge\bfseries}[\vspace{4pt}{\color{accent}\titlerule[1.2pt]}]
\titlespacing*{\chapter}{0pt}{-10pt}{22pt}
\titleformat{\section}{\sffamily\Large\bfseries\color{navy}}{\color{accent}\thesection}{0.7em}{}
\titleformat{\subsection}{\sffamily\large\bfseries\color{navy}}{\color{accent}\thesubsection}{0.7em}{}
\titleformat{\subsubsection}{\sffamily\normalsize\bfseries\color{navy}}{\thesubsubsection}{0.6em}{}
\titlespacing*{\section}{0pt}{16pt}{7pt}
\titlespacing*{\subsection}{0pt}{12pt}{5pt}
\setcounter{secnumdepth}{3}
\setcounter{tocdepth}{2}

% --- колонтитули -----------------------------------------------------------
\usepackage{fancyhdr}
\pagestyle{fancy}
\fancyhf{}
\renewcommand{\headrulewidth}{0.4pt}
\renewcommand{\headrule}{{\color{rule}\hrule width\headwidth height\headrulewidth}}
\fancyhead[L]{\sffamily\footnotesize\color{muted}ML откриване на прониквания в трафика на център за данни}
\fancyhead[R]{\sffamily\footnotesize\color{muted}\nouppercase{\leftmark}}
\fancyfoot[C]{\sffamily\footnotesize\color{navy}\thepage}
\fancypagestyle{plain}{\fancyhf{}\renewcommand{\headrulewidth}{0pt}\fancyfoot[C]{\sffamily\footnotesize\color{navy}\thepage}}
\renewcommand{\chaptermark}[1]{\markboth{\thechapter.\ #1}{}}

% --- кутии -----------------------------------------------------------------
\usepackage[most]{tcolorbox}
\newtcolorbox{keybox}[1][Основни изводи]{enhanced,breakable,colback=boxbg,colframe=accent,boxrule=0pt,
  leftrule=3pt,arc=1pt,left=8pt,right=8pt,top=5pt,bottom=5pt,fonttitle=\sffamily\bfseries\small,
  coltitle=navy,colbacktitle=boxbg,title=#1,attach title to upper={\par\smallskip},fontupper=\small}
\newtcolorbox{notebox}[1][Бележка]{enhanced,breakable,colback=white,colframe=orange,boxrule=0.6pt,arc=1pt,
  left=8pt,right=8pt,top=4pt,bottom=4pt,fonttitle=\sffamily\bfseries\small,coltitle=orange,
  colbacktitle=white,title=#1,attach title to upper={\par\smallskip},fontupper=\small}
\newtcolorbox{interp}{enhanced,colback=white,colframe=white,boxrule=0pt,borderline west={2pt}{0pt}{rule},
  left=10pt,right=4pt,top=2pt,bottom=2pt,fontupper=\small}
\newcommand{\interpret}[1]{\begin{interp}\textsf{\textbf{\color{navy}Какво показва фигурата.}}\ #1\end{interp}}
\newcommand{\interprettab}[1]{\begin{interp}\textsf{\textbf{\color{navy}Как се чете таблицата.}}\ #1\end{interp}}

% --- код -------------------------------------------------------------------
\usepackage{listings}
\lstdefinestyle{py}{language=Python,basicstyle=\ttfamily\footnotesize,keywordstyle=\color{accent}\bfseries,
  commentstyle=\color{muted},stringstyle=\color{orange},showstringspaces=false,frame=l,
  framerule=2pt,rulecolor=\color{accent},backgroundcolor=\color{boxbg},xleftmargin=12pt,framexleftmargin=8pt,
  breaklines=true,columns=fullflexible,extendedchars=true}
\lstset{style=py}

% --- таблици ---------------------------------------------------------------
\newcolumntype{R}{>{\raggedleft\arraybackslash}X}
\newcolumntype{L}{>{\raggedright\arraybackslash}X}
\newcommand{\hd}[1]{\makecell{\sffamily\bfseries\color{white}#1}}
\newcommand{\headrow}{\rowcolor{navy}}
\renewcommand{\arraystretch}{1.18}

% --- числа от резултатите (генерирани автоматично) ------------------------
\newcommand{\val}[1]{\ifcsname val@#1\endcsname\csname val@#1\endcsname\else\textbf{\color{red}??#1??}\fi}

% --- математика ------------------------------------------------------------
\newcommand{\R}{\mathbb{R}}
\newcommand{\E}{\mathbb{E}}
\newcommand{\vect}[1]{\mathbf{#1}}
\DeclareMathOperator*{\argmax}{arg\,max}
\DeclareMathOperator*{\argmin}{arg\,min}
\DeclareMathOperator{\sign}{sign}
\newcommand{\txt}[1]{\text{\normalfont #1}}

\usepackage[hidelinks]{hyperref}
\hypersetup{colorlinks=true,linkcolor=navy,citecolor=accent,urlcolor=accent,
  pdftitle={Експериментално изследване на откриване на проникване в мрежов трафик на център за данни, базирано на машинно обучение},
  pdfauthor={Kristiyan Yanutsev}}
\usepackage[english,nameinlink,noabbrev]{cleveref}
\crefname{equation}{формула}{формули}
\crefname{figure}{фиг.}{фиг.}
\crefname{table}{табл.}{табл.}
\crefname{section}{раздел}{раздели}
\crefname{chapter}{глава}{глави}
\crefname{appendix}{приложение}{приложения}
\creflabelformat{equation}{(#2#1#3)}
\pretocmd{\section}{\Needspace{20\baselineskip}}{}{}\pretocmd{\subsection}{\Needspace{12\baselineskip}}{}{}

% Без курсив никъде в текста
\renewcommand{\emph}[1]{#1}
\renewcommand{\textit}[1]{#1}
\renewcommand{\textsl}[1]{#1}
\AtBeginDocument{\renewcommand{\itshape}{\upshape}\renewcommand{\slshape}{\upshape}}
